\documentclass[../main.tex]{subfiles}

\begin{document}
\begin{problema}
    Sea \(u \colon \Omega \subseteq \mathbb{C} \to \mathbb{R}\) armónica en \(\Omega\) y
    \(z_{0}\in \Omega\).

    \begin{enumerate}
        \item Pruebe que, si \(0 < r\) es tal que \(\overline{B_{r}(z_{0})} \subseteq
            \Omega\), entonces

            \begin{equation*}
                u(z_{0}) = \dfrac{1}{2\pi}\int_{0}^{2\pi}u \bigl(r \mathrm{e}^{it} +
                z_{0}\bigr)\odif{t}.
            \end{equation*}

        \item Pruebe que, si \(u(z) \leq u(z_{0})\) para toda \(z\in
            \Omega\), entonces \(u\) es
            constante en \(\Omega\).
    \end{enumerate}
\end{problema}

\startsolution

\begin{enumerate}
    \item Como \(\Omega\) es simplemente conexo, \(u = \Re f\) para una función \(f
        \colon \Omega \to \mathbb{C}\) analítica y aplicando la fórmula
        integral de Cauchy,

        \begin{equation*}
            u(z_{0}) = \dfrac{1}{2\pi}\int_{0}^{2\pi}u(z_{0} + r
            \mathrm{e}^{it})\odif{t}.
        \end{equation*}

    \item Consideremos un disco abierto \(B_{r}(z_{0}) \subseteq \Omega\). Como \(B\) es
        simplemente conexo, entonces existe \(f\) analítica en
        \(B_{r}(z_{0})\) tal que \(u = \Re f\).
        Entonces la función \(\mathrm{e}^{f(z)}\) también es analítica y
        \(\abs{\mathrm{e}^{f(z)}} = \mathrm{e}^{u(z)}\). Como \(u\) alcanza un máximo en
        \(z_{0}\) y la función exponencial real es creciente,
        \(\abs{\mathrm{e}^{f(z)}}\)
        también alcanza un máximo en \(z_{0}\). Por el principio del módulo máximo,
        \(\mathrm{e}^{f(z)}\) es constante en \(B_{r}(z_{0})\) y por lo tanto
        \(\mathrm{e}^{u(z)}\)
        también lo es, por lo que \(u\) es constante en \(B_{r}(z_{0})\).
        Podemos aplicar este
        argumento a cualquier punto del conjunto

        \begin{equation*}
            A = \bigl\lbrace z\in \Omega \mid u(z) = u(z_{0})\bigr\rbrace
        \end{equation*}

        lo que implica que \(A\) es abierto en \(\Omega\). Pero claramente
        \(A\) también es
        cerrado y no vacío. Como \(\Omega\) es conexo, \(A = \Omega\)  y
        \(u\) es constante.
\end{enumerate}
\end{document}
